\documentclass{article}
\usepackage{iclr2026_conference}   % swap for the OFFICIAL ICLR 2026 .sty at submission
\usepackage{booktabs}
\usepackage{amsmath,amssymb}
\usepackage{graphicx}
\usepackage[numbers]{natbib}
\usepackage{hyperref}
\usepackage{xcolor}
\newcommand{\numcells}{18}
\newcommand{\nummodels}{6}
\newcommand{\medsave}{50\%}
\newcommand{\numtraps}{5}
\newcommand{\maxpricex}{8.8$\times$}
\newcommand{\maxlatx}{68$\times$}
                      % \medsave \numcells \nummodels \numtraps \maxpricex \maxlatx

\title{Market-Aware LLM Inference Routing:\\ You Cannot Pick a Provider From the Price List}
\author{Tony Wen \\ \texttt{ts2015656@gmail.com}}

\begin{document}
\maketitle

\begin{abstract}
Existing LLM routers treat each model's cost as a static per-model attribute and select a
model by query difficulty. We argue this misses a first-order effect of the emerging,
commoditized compute market: for open-weight models served by \emph{competing providers},
the same model varies sharply in price, latency, reliability, and---silently---in answer
\emph{quality}. We formalize \emph{market-aware routing} as a price-\emph{taker} decision
problem (the router optimizes against exogenous, provider-given prices, distinguishing it
from provider-side pricing games) and show, on \numcells{} (model$\times$task) cells across
\nummodels{} open models and competing providers measured live through a public aggregator,
that price, latency, and quality are mutually uncorrelated: cross-provider price spreads reach
\maxpricex{}, latency spreads \maxlatx{}, and ``the same model'' ranges from 35\% to 90\%
accuracy on hard reasoning. Consequently neither price nor the advertised quantization label
predicts quality, and ``route to the cheapest'' silently drops below a quality floor in
\numtraps{} of our cells. A measurement-driven router that picks the cheapest provider which is
\emph{both} measured-equivalent \emph{and} currently healthy realizes a median \medsave{} cost
reduction at matched quality. We release the measurement harness, the consolidated map, and an
OpenAI-compatible routing gateway. We are explicit that our evidence is cross-sectional, not
temporal; we discuss what the (not-yet-live) token spot market would add.
\end{abstract}

\section{Introduction}
LLM inference is being commoditized. GPU rental prices are published as daily indices and a
GPU \emph{compute futures} market was announced in 2026 \citep{cme2026compute}; the same
open-weight model (Llama, Qwen, DeepSeek) is now served by dozens of competing providers at
prices that differ by up to an order of magnitude. Yet the dominant line of LLM routing work
---FrugalGPT \citep{chen2023frugalgpt}, RouteLLM \citep{ong2024routellm}, Hybrid~LLM
\citep{ding2024hybrid}, CARROT \citep{somerstep2025carrot}, MixLLM \citep{wang2025mixllm}---
represents cost as a \emph{static} per-model quantity (a published price, or the fraction of
traffic sent to an expensive tier) and routes on query difficulty alone.

This paper makes a simple empirical claim with a sharp consequence: \textbf{across competing
providers, price, latency, reliability, and quality of the same model are mutually
uncorrelated, so a provider cannot be chosen from the price list---it must be measured.} The
cheapest endpoint is typically the most aggressively quantized (fp8/fp4) and sometimes silently
loses accuracy; but the most \emph{expensive} endpoint is often neither the best nor reliable.
Only direct, per-task measurement separates the safe-cheap provider from the cheap-and-degraded
one.

\paragraph{Contributions.}
(i) We formalize market-aware routing as a \emph{price-taker} constrained decision problem
(Section~\ref{sec:form}), with a \emph{price-taker assumption} that cleanly separates it from
provider-side pricing games \citep{prillm2025} and from infrastructure-layer spot serving
\citep{miao2024spotserve,liu2025skyserve}. (ii) We measure \numcells{} (model$\times$task) cells
across \nummodels{} open models and their competing providers live (Section~\ref{sec:meas}),
and show price$\perp$quality$\perp$latency$\perp$reliability (Section~\ref{sec:findings}).
(iii) We build a measurement-driven router and show a median \medsave{} cost reduction at
matched measured quality while avoiding the \numtraps{} ``cheapest-is-degraded'' traps
(Section~\ref{sec:router}). We release the harness, the map, and an OpenAI-compatible gateway.

\section{Related Work}
\paragraph{Routing under static cost.} The client-side routing lineage selects a model per
query under a static cost model: a published per-token price or the offload fraction to a
cheaper tier \citep{chen2023frugalgpt,ong2024routellm,ding2024hybrid,somerstep2025carrot,
wang2025mixllm}. None model time-varying prices, provider capacity, or congestion-dependent
latency, and---central to our finding---none model \emph{cross-provider} quality heterogeneity
for a fixed model. CARROT's SPROUT benchmark is ``price-aware'' only in the sense of folding
static list prices into the estimate; this is orthogonal to measuring a live market.

\paragraph{Provider-side pricing.} A recent line studies how a \emph{provider} should price its
service. PriLLM \citep{prillm2025} models the market as a Stackelberg game where the provider is
the leader \emph{setting} price to maximize profit. This is the \emph{inverse} of our problem:
we are a client/router that takes the market price as exogenous input. Price is the leader's
decision variable there; it is our environment's state here.

\paragraph{Spot/preemptible serving.} SpotServe \citep{miao2024spotserve} and SkyServe
\citep{liu2025skyserve} cut cost by running inference on cheap preemptible GPUs and adapting to
\emph{availability}, assuming a \emph{static} spot-vs-on-demand gap. They act on
\emph{where/how} a model is served, not \emph{which} provider answers a query, and sit on the
opposite side of the routing boundary. Reacting to exogenous time-varying prices is also
established in carbon-/electricity-aware datacenter scheduling, but that acts on
operator-controlled physical resources, not client-side cross-provider API routing.

\section{Problem Formulation}\label{sec:form}
Queries $x_1,\dots,x_T$ arrive over a horizon. A pool $M=\{m_1,\dots,m_K\}$ of model/provider
endpoints each has a market state $(p_m(t),c_m(t),\ell_m(t),\phi_m(t))$: token price, available
capacity, congestion-dependent latency, failure rate. Each query carries a task type, a quality
requirement $\theta$, a deadline $d$, and a value $v$.

\paragraph{Assumption 1 (price-taker / exogenous market).} The market state evolves
independently of the router's action: $P(m(t{+}1)\mid m(t),a_t)=P(m(t{+}1)\mid m(t))$. A single
client's dispatch does not move provider prices or aggregate congestion. This one line is the
formal separation from provider-side pricing \citep{prillm2025}, where price is endogenous.

\paragraph{Objective.} With state $s_t=(x_t,m(t),b_t,\rho_t)$ (query, market, remaining budget,
reserved inventory) and action $a_t$ (provider, spot/reserved tier, defer, fallback), we
maximize value-weighted quality subject to budget, SLO, and quality-floor constraints:
\[
\max_\pi\ \mathbb{E}\Big[\textstyle\sum_t v_t\,\mathrm{score}(y_{a_t},y^\star)\Big]
\ \ \text{s.t.}\ \ \mathbb{E}\big[\textstyle\sum_t \mathrm{cost}\big]\le B,\ \
\Pr(\ell_{a_t}\!>d_t)\le\delta,\ \ Q(x_t,a_t)\ge\theta_t.
\]

\paragraph{Proposition 1 (bandit reduction).} If there is no shared budget, no reserved
inventory, the market is exogenous (Assumption 1), and no deferral, the constrained MDP
decouples into independent per-query problems and the optimal policy is the greedy per-query
utility maximizer $a_t^\star=\arg\max_m U_m(x_t,m(t))$; a contextual bandit over
$(x_t,m(t))$ suffices. Shared budget, reserved inventory, or deferral break this and require a
constrained policy. Our empirical study operates in the bandit regime; we treat the per-query
utility as $Q_m$ measured offline minus price and latency penalties at the current market state.

\paragraph{When does measurement matter?} Let $\pi_{\text{static}}$ route on the time-averaged
price and $\pi_{\text{mkt}}$ on the realized state. The advantage of measurement is zero only
when the arg-max provider is unchanged by the realization, and otherwise grows with the
\emph{dispersion} of price/quality/latency across providers and the frequency with which the
optimal provider flips. The next sections show that dispersion is large \emph{right now}, even
before any temporal price dynamics.

\section{Measurement Methodology}\label{sec:meas}
We probe each endpoint through a public multi-provider aggregator (OpenRouter
\citep{openrouter}), pinning each request to a single provider with fallbacks disabled, so a
measurement reflects \emph{that} provider's serving stack (quantization, kernels, sampling).
For each (model, task) we run an objective probe set ($n{=}15$--$20$, exact-match scored) over
all competing providers and record accuracy, mean latency, real per-call cost, and availability
(served fraction, with retry/back-off and error typing: 429/5xx/timeout/client). Tasks span a
difficulty range: multi-step \emph{math} (reasoning), \emph{extraction}, and \emph{sentiment
classification}. Endpoints are deduplicated by (provider, quantization, price).

\paragraph{Honest measurement caveats.} (a) Probe sets are small ($n\le20$): coarse
distinctions (e.g.\ 50\% vs.\ 55\%) are within binomial noise, while the large gaps we
emphasize (35\% vs.\ 90\%) are several standard errors. (b) Firing $\sim$20 rapid pinned
requests at one provider can self-induce rate limits; we early-abort a provider after three
consecutive failures and report it as low-availability, which may understate true uptime.
(c) Reasoning models require a sufficient token budget; an initial truncation bug collapsed a
strong model to 15\% before we corrected it---an instructive reminder that measurement harnesses
are themselves a source of error.

\section{Empirical Findings}\label{sec:findings}
Table~\ref{tab:main} reports all \numcells{} cells. Three patterns hold.

\begin{table}[t]\centering\small
\caption{Measured dispersion per (model$\times$task). $|P|$ = competing providers;
Acc.\ range = min--max accuracy of \emph{the same model}; Price$\times$/Lat.$\times$ =
max/min spread; Route price = price of the cheapest measured-equivalent healthy provider;
Save = its reduction vs.\ the premium (most expensive) provider at matched quality.
$\dagger$ marks cells where the cheapest provider falls below the quality floor.}
\label{tab:main}
\begin{tabular}{llrccccr}
\toprule
Model & Task & $|P|$ & Acc.\ range & Price$\times$ & Lat.$\times$ & Route price & Save \\
\midrule
deepseek-chat-v3.1 & classification & 7 & 100\%--100\% & 2.3 & 68 & \$0.500 & 57\% \\
 & extraction & 6 & 93\%--100\% & 1.8 & 34 & \$0.625 & 46\% \\
 & math & 7 & 100\%--100\% & 2.3 & 24 & \$0.625 & 46\% \\
gemma-3-27b-it & classification & 5 & 100\%--100\% & 2.5 & 4 & \$0.120 & 61\% \\
 & extraction & 5 & 100\%--100\% & 2.5 & 8 & \$0.120 & 61\% \\
 & math & 4 & 53\%--72\%\,$\dagger$ & 1.9 & 3 & \$0.305 & 0\% \\
llama-3.1-8b-instruct & classification & 5 & 100\%--100\% & 8.8 & 8 & \$0.025 & 89\% \\
 & extraction & 5 & 80\%--93\%\,$\dagger$ & 8.8 & 10 & \$0.035 & 84\% \\
 & math & 5 & 56\%--85\%\,$\dagger$ & 8.8 & 13 & \$0.220 & 0\% \\
llama-3.3-70b-instruct & classification & 9 & 100\%--100\% & 4.8 & 25 & \$0.265 & 79\% \\
 & extraction & 10 & 93\%--100\% & 5.0 & 6 & \$0.210 & 80\% \\
 & math & 11 & 40\%--75\%\,$\dagger$ & 6.1 & 56 & \$0.265 & 79\% \\
llama-4-maverick & classification & 4 & 100\%--100\% & 2.0 & 3 & \$0.375 & 50\% \\
 & extraction & 4 & 93\%--93\% & 2.0 & 4 & \$0.375 & 50\% \\
 & math & 4 & 80\%--85\% & 2.0 & 2 & \$0.375 & 50\% \\
mistral-small-3.2-24b-instruct & classification & 4 & 100\%--100\% & 1.5 & 2 & \$0.172 & 14\% \\
 & extraction & 4 & 93\%--100\%\,$\dagger$ & 1.5 & 3 & \$0.200 & 0\% \\
 & math & 4 & 100\%--100\% & 1.5 & 4 & \$0.138 & 31\% \\
\bottomrule
\end{tabular}
\end{table}

\paragraph{(1) Quality dispersion is selective but large.} On classification, providers
saturate (all $\approx$100\%) even for a weak 8B model. On hard multi-step math, ``the same
model'' ranges 35--90\% (8B), 40--75\% (70B), 53--72\% (others): aggressive quantization and
serving choices silently cost tens of points. Dispersion is driven by \emph{task difficulty},
not model size alone---the headline 35--90\% gap is a weak-model$\times$hard-task worst case.

\paragraph{(2) Price and latency dispersion are universal.} Even on saturated tasks, price
spreads reach \maxpricex{} and latency spreads \maxlatx{} for the identical model. The most
expensive provider is frequently neither the most accurate nor the most reliable.

\paragraph{(3) Neither price nor quant label predicts quality.} Two providers serving the same
``fp8'' model differ by 30 accuracy points; a bf16 endpoint can be the most expensive and among
the least accurate. Only direct measurement orders providers correctly.

\section{Market-Aware Router}\label{sec:router}
We route each request to the cheapest provider whose measured accuracy is within a floor of the
best \emph{and} whose availability exceeds a threshold. Against a ``premium'' baseline (pay up
for the dearest provider) the router achieves a median \medsave{} cost reduction at matched
measured quality; against a ``price-blind'' baseline (sort by price) it avoids the \numtraps{}
cells where the cheapest provider is silently sub-floor. The same policy, exposed as an
OpenAI-compatible endpoint, lets a caller change one line (\texttt{base\_url}) and tag the task
to obtain the per-task floor; a savings ledger records the counterfactual cost. On capable
models where quality saturates, the router instead trades a little price for reliability,
avoiding cheap-but-rate-limited endpoints.

\section{Limitations}
Our evidence is \emph{cross-sectional}: it establishes large dispersion \emph{at an instant},
not that the cheapest-safe provider \emph{churns over time}. The latter---the temporal,
``dynamic price'' story motivated by compute futures \citep{cme2026compute}---requires a
longitudinal trace we have not yet collected, and a live token spot market does not yet exist
(token list prices change on the order of days). Our measurements are also point-in-time and
provider availability is partly confounded by self-induced rate limits. We therefore frame the
contribution as measurement-driven routing under \emph{static-but-heterogeneous} markets, with
the dynamic regime as motivated future work.

\section{Conclusion}
The same open model, served by competing providers, differs enough in price, latency,
reliability, and silent quality that a provider cannot be chosen from the price list. Measuring
the (model$\times$task$\times$provider) cell and routing to the cheapest measured-equivalent
healthy endpoint yields a median \medsave{} saving at matched quality. The measurement map is
the durable asset; the router and gateway are thin layers on top.

\bibliographystyle{plainnat}
\bibliography{references}
\end{document}
